\Section{Discussion and Conclusion}
\label{sec:conclusion}

\SubSection{Evaluation and Open Questions}

We have applied the approach to examples from the basic Move libraries, covering loops, higher-order functions, arithmetic, state-dependent algorithms, and more.
While the \WP algorithm is stable and well-tested, the definition of skills and instructions is an ever-moving target.

\Paragraph{Ordering of Inference Tasks}
Currently the agent walks through the inference tasks --- loop invariants, \WP analysis, simplification, and verification --- in a fixed order.
We arrived at this order experimentally. In one experiment, \WP first derived basic specs, then the model inferred invariants, and \WP ran again.
The vacuous specifications from the first run were unhelpful and led the model into unfavourable chains of thought.

As workflows of this kind become more common and foundation models ingest them, we expect to leave ordering to the model's planning abilities rather than stack heuristics on heuristics.
Technically, each step is transactional, so the agent can execute them in any order. We leave this to future research.

\Paragraph{Effectiveness of the Hybrid Approach}
A conceptual question is to what extent the agent actually uses \WP-generated specs as a baseline, and whether \WP is necessary at all --- could the agent do all the work, as in~\cite{PeiEtAl23, KamathEtAl23, MSG25}?
We observe that the \WP baseline keeps the agent from suggesting outright wrong specifications, but we lack the data to answer conclusively.
A systematic evaluation requires a carefully designed case study, addressing questions such as:
\begin{Itemize}
\item What is the success rate of the hybrid approach compared to pure LLM inference?
\item Does the token cost improve when \WP is involved? Or does it get worse because of the additional context?
\item Is the set of samples representative? Is the model biased toward the given samples because it already knows the solution (e.g. for standard problems like the |find| function in Sec.~\ref{sec:find})?
\item Is the next generation of models changing the picture? The work here has been done with Opus 4.7; what will happen when Opus 5.x hits the market later in 2026?
\item How does the approach scale for larger systems? How does it compare to handwritten specs?
\end{Itemize}

\noindent We plan to apply inference retrospectively to the Aptos Framework \cite{AptosFrameworkBook}, which already has specifications, and compare against the existing ones.
We also plan to apply it to a newer in-house project, a complex perpetual trading engine, with results to be reported in future publications.

%  but one example from a small controlled study suggests the potential usefulness of the hybrid approach.
% Both conditions ran as isolated Claude Sonnet~4.6 sessions with identical documentation access; the sole difference was that the hybrid condition (\textbf{H}) could invoke the \WP MCP tool and the agent-only condition (\textbf{A}) could not.
%
% Consider !fold!, which left-folds a !vector<u64>! with a closure !f: |u64,u64|u64! that may abort.
% A complete abort contract requires a safety loop invariant certifying that no prefix application of !f! has aborted:
% \begin{MoveBox}
% invariant forall j: u64 where j < i:
%     !aborts_of<f>(spec_fold(v, f, init, j), v[j]);
% \end{MoveBox}
% It is derived by simplifying the one that \WP derives mechanically by backward analysis from the abort postcondition.
% \textbf{H} received it directly and verified in two iterations (each iteration is one edit-then-verify cycle against the prover); \textbf{A} independently discovered the !spec_fold! helper and the existential abort condition but needed one additional iteration to arrive at the same safety invariant from a failing proof.

% Although more case studies are needed to draw firm conclusions, this example illustrates where the hybrid approach has potential: \WP pre-fills structural invariants that follow mechanically from the call graph but are non-obvious to discover from scratch.


\Paragraph{Skill Design}
Writing effective skills is something of a `black art'. A few lessons have emerged.

\begin{Itemize}
\item \emph{Minimize context.}
The skill should carry only the information the agent needs --- not a tour of the Move ecosystem.
We repeatedly found that extra background dilutes attention, lengthens the working context, and gives the agent room to diverge from the task.

\item \emph{Task first, then supporting material bottom-up.}
The skill opens with a concise statement of what to do --- the task list of Sec.~\ref{sec:skills} --- with forward references to the rules and reference material that follow.
The body is then ordered bottom-up: spec-language reference first, then tool documentation, then domain-specific guidelines, so that by the time the agent revisits the task list every term has been introduced. With top-down organization, the agent tends to draw conclusions and act on them before reading all the relevant information.

\item \emph{Restate the prohibitions.}
Agents tend to take shortcuts when searching for a solution.
The agent often tries to bypass a hard verification problem by weakening a condition --- silencing an abort case, narrowing a postcondition --- so the prover reports success.
We therefore restate this prohibition at every relevant point in the skill text.
Even so, we still occasionally catch \CC taking the shortcut; the user-visible task tracker and the verification oracle remain the last line of defense.

\item \emph{Testing.}
Automated testing of the skill system is hard: the agent's choices are non-deterministic and coding-agent usage is metered, making CI integration impractical. Our best lever for stabilization is extended manual application.
\end{Itemize}

% We currently have a project under way that re-specifies larger parts of the Move and Aptos frameworks; continuous human feedback from this activity flows into the evolution of the plugin.
%

\SubSection{Related Work}

Specification inference has a long history.
Daikon~\cite{DaikonTSE01} pioneered \emph{dynamic} invariant detection, learning likely invariants from program traces; subsequent work explored static inference, abstract interpretation, and template-based synthesis, but no unified approach has emerged --- tools have relied on heuristics tuned to particular property classes and source languages.

The advent of large language models has revived interest in the area.
Pei et al.~\cite{PeiEtAl23} showed that fine-tuned LLMs can predict program invariants statically at quality competitive with dynamic analysis.
Kamath et al.~\cite{KamathEtAl23} pair an LLM that \emph{proposes} loop invariants with a symbolic checker that \emph{validates} them, retrying on failure --- a propose-and-validate loop that recurs in much of the subsequent work.
For Verus, AlphaVerus and VeriStruct~\cite{VeriStruct26} synthesize Verus-style specifications (and proofs) for Rust functions, modules, and data structures, using fine-tuned LLMs, tree search, or planner-orchestrated phases.
For smart contracts in Solidity, PropertyGPT~\cite{PropertyGPT25} retrieves human-written Certora properties as input examples for specification generation.
Most directly related to our work, MSG~\cite{MSG25} targets the same Move specification language through a skill-based agentic design that consumes verifier feedback beyond a binary success/failure verdict.
% Unlike the hybrid approach we pursue, MSG leaves all specification generation to the LLM; it does not exploit that --- once loop invariants are in place --- weakest-precondition computation is mechanical and well-defined, and yields a substantial fragment of the specification (abort conditions, modifies clauses, consequence-style postconditions) at essentially no token cost and with full precision.

Across these systems the LLM remains responsible for the entire specification; the verifier serves as a binary oracle or, in MSG and our setup, as a feedback channel, but never as a contributor of mechanically derived content.

A separate line of work, exemplified by DafnyPro~\cite{DafnyPro26} and AutoVerus~\cite{AutoVerus25}, tackles the complementary problem of synthesizing the auxiliary proof annotations a verifier needs when the specification is already given. Our plugin can also synthesize proof hints (Sec.~\ref{sec:pow}), but this is not the focus of this paper.

Reference elimination, the bytecode transformation underlying our \WP analysis (Sec.~\ref{sec:wp}), was introduced in~\cite{DillGPQXZ22} and is conceptually shared with Aeneas~\cite{Aeneas}, which compiles Rust to a pure functional model in a target prover (Lean, F*, Coq, HOL4) via forward and backward functions.
Aeneas does \emph{not} infer specifications: the user writes them in the target prover, against the translated functional model. We expect our \WP approach can be applied similarly to Rust.

Our backend IVL, Boogie~\cite{boogie,boogieIVL}, also performs \WP on the procedural code passed in~\cite{BL05} and passes the result on to the SMT solver. One major difference from our approach is that we aim to generate user-readable conditions, in terms of Move, whereas Boogie's \WP is not intended for external consumption. This is reflected in our mapping of references back to the source language, and in the significant effort put into expression simplification, since raw \WP predicates can be very verbose.

In summary, two aspects distinguish our setup from the work above.
First, to our knowledge, no prior LLM-based spec-inference system uses a sound mechanical baseline alongside the LLM.
\WP itself produces a substantial part of the inferred specification --- abort conditions, modifies clauses, and consequence-style postconditions --- so the AI is left only with what \WP cannot recover (loop invariants, high-level conservation laws), and the two streams are validated together by the Move Prover.
The AI side runs a propose-observe-refine loop like the systems above; the mechanical baseline narrows the LLM's search space.

Second, the workflow is delivered as a plugin to a general-purpose, interactive coding agent (Claude Code) via skills, edit-time hooks, and an MCP server, rather than as a custom-built tool: the agent participates in the loop and -- most importantly -- the user can intervene at any point.
To our knowledge, no prior LLM-based inference system for verification languages has this kind of coding-agent integration, including agentic ones such as MSG.

\SubSection{Conclusion}

We have presented early work on a specification inference tool for the Move Prover that combines a sound, mechanical \WP analysis over Move bytecode with an agentic coding CLI; the prover serves as the oracle for both signals.
The two inference streams compose in a single instrumentation pipeline, and the workflow is delivered as a Claude Code plugin --- skills, hooks, and an MCP server --- so inference runs inside the same conversational loop the user already writes Move in.
We have applied the system to canonical Move libraries spanning loops, higher-order functions, arithmetic, references, and global state. Future work will stabilize and refine the tool by applying it to the large corpus of safety-relevant Move code used at Aptos for perpetual trading and onchain money management.

%%% Local Variables:
%%% mode: latex
%%% TeX-master: "main"
%%% End:
